\chapter{Research Onboarding and Question Design}

\WeekHeader{1}{Research Onboarding and Question Design}{Turn a broad interest into a testable chemical research question.}{One-page project brief with question, system, candidate chemicals, expected evidence, and risks.}

\section{Why This Matters}
Strong computational research starts before any code is run. Students need a clear scientific question, a defensible reason for studying a chemical set, and a plan for evaluating evidence. Without this foundation, database enrichment becomes a list of facts rather than a research workflow.

\section{Concepts}
\begin{itemize}
  \item \textbf{Research system}: the organism, environment, treatment, material, food, exposure, sample type, or process being studied.
  \item \textbf{Chemical evidence}: tentative identifications, mass spectral matches, database annotations, literature support, functional groups, biological activities, hazard terms, occurrence records, pathways, reactions, and measured properties.
  \item \textbf{Analysis-ready variable}: a clean field that can be grouped, counted, filtered, modeled, or visualized.
  \item \textbf{Provenance}: the source and retrieval context for a piece of evidence.
  \item \textbf{Uncertainty}: the gap between a database hit and a confirmed scientific conclusion.
\end{itemize}

\section{uafR and R Skills}
Students do not need to run the full workflow in Week 1. They should learn what the tool can produce:
\begin{itemize}
  \item Core GC-MS workflow: \texttt{spreadOut()}, \texttt{mzExacto()}, and \texttt{exactoThese()}.
  \item Chemical enrichment: \texttt{categorate(detail = "research")}.
  \item Research trait tables: \texttt{ChemicalTraits}, \texttt{ChemicalTraitOntology}, and \texttt{ChemicalTraitMatrix}.
  \item Hazard summary table: \texttt{ChemicalHazards}.
  \item Occurrence summary table: \texttt{ChemicalOccurrences}.
  \item Measurement summary table: \texttt{ChemicalMeasurementSummary}.
  \item KEGG-specific interpretation table: \texttt{KEGGReactionParticipants}.
  \item Quality checks: \texttt{validateCategorateResult()} and \texttt{TableQuality}.
\end{itemize}

\section{Guided Lab}
\begin{enumerate}
  \item Choose one research domain: plant metabolites, environmental exposure, food chemistry, microbial products, pharmacology, materials, forensic chemistry, toxicology, or another reviewed topic.
  \item Write three possible research questions. Each question must name the system, chemicals or chemical class, comparison, and expected evidence.
  \item Identify a starting chemical list. This can come from GC-MS results, literature, standards, known metabolites, or a small exploratory list.
  \item Mark each chemical as known, tentative, uncertain, or exploratory.
  \item Write a risk note: what could make the analysis misleading?
\end{enumerate}

\begin{goldbox}{Example Research Questions}
\begin{itemize}
  \item Do volatile compounds detected in two plant treatments differ in sensory trait profiles, natural product occurrence, or pathway roles?
  \item Which detected compounds have safety or exposure signals that require additional review before biological interpretation?
  \item Do compounds associated with a sample group cluster by functional group, pathway, or literature-backed biological activity?
\end{itemize}
\end{goldbox}

\section{Independent Extension}
Students interview the program lead, the dsDNA Core mentor, or a project lead for 20 minutes and capture:
\begin{itemize}
  \item The research context.
  \item Why the chemical list matters.
  \item What decisions the analysis should support.
  \item What evidence would be convincing.
  \item What evidence would be weak or unacceptable.
\end{itemize}

\section{Deliverables}
\begin{tabularx}{\textwidth}{p{0.25\textwidth}YY}
\DeliverableTableHeader
Project brief & Includes question, system, comparison, candidate chemical list, and expected output. & PDF or Markdown\\
Chemical intake table & Has one row per chemical and columns for source, confidence, project relevance, and notes. & CSV\\
Risk note & Names at least three ways the analysis could be wrong or overinterpreted. & Research log\\
\end{tabularx}

\section{Quality Checklist}
\begin{checklistbox}{Before Week 2}
\begin{itemize}
  \item The question is narrow enough to analyze in 10 weeks.
  \item The chemical list has a clear source and inclusion rule.
  \item The student can explain why chemical traits, database evidence, or GC-MS patterns matter for the question.
  \item The project scope has passed dsDNA Core review.
\end{itemize}
\end{checklistbox}
